The polynomial in \(\mathrm{def:first\text{-}bin\text{-}alias}\) is
\[
 \phi(u)=1-3u+2u^2,
 \qquad
 \int_0^1u\phi(u)\,du
 =\int_0^1(u-3u^2+2u^3)\,du
 =\frac12-1+\frac12=0.
\]
Also \(\phi(1)=0\), \(|\phi(u)|\leq1\) on \([0,1]\), and \(\phi'(u)=-3+4u\), so \(\sup_{u\in[0,1]}|\phi'(u)|=3\).  Define
\[
 a_\Delta(r)=\Delta\phi(r/\Delta)\mathbf1\{0\leq r<\Delta\}.
\]
The two pieces meet at \(r=\Delta\) because \(\phi(1)=0\); hence the derivative bound shows that \(a_\Delta\) is continuous and \(3\)-Lipschitz on \([0,1]\).  The reverse triangle inequality gives
\[
 \bigl|\lVert x-b\rVert_2-\lVert x'-b\rVert_2\bigr|
 \leq\lVert x-x'\rVert_2.
\]

Consider either sign in \(\mathrm{def:first\text{-}bin\text{-}alias}\).  Choose \(\theta_0=1/2\) and
\[
 \theta_{1,\pm}=\frac12\pm c_L\Delta
 =\frac12\pm c_La_\Delta(0).
\]
For every treated point \(x\) of radius \(r\leq1/2\),
\[
 |\mu_{1,\pm}(x)-\theta_{1,\pm}|
 =c_L|a_\Delta(r)-a_\Delta(0)|
 \leq3c_Lr
 \leq Lr,
\]
where \(3c_L\leq3L/8\leq L\).  The control deviation is zero.  Furthermore,
\[
 |c_La_\Delta(r)|\leq c_L\Delta
 \leq\frac1{16}\frac14=\frac1{64},
\]
so the constructed Bernoulli means lie in \([31/64,33/64]\subset[0,1]\).  Thus the bounded-outcome conditional-mean and point-anchor conditions hold.

The construction supplies the known half-plane partition, sharp assignment, consistency, and conditionally Bernoulli potential outcomes.  Uniformity of \(X\) on the unit disc and polar integration give, for \(t\in\{0,1\}\) and \(0<H\leq1/2\),
\[
 \mathbb P(T=t,R<H)
 =\frac1\pi\int_{\Theta_t}\int_0^H r\,dr\,d\vartheta
 =\frac{H^2}{2}
 \geq c_0H^2,
 \qquad |\Theta_t|=\pi,
\]
because \(c_0\leq1/2\).  Hence both constructed laws belong to \(\mathcal P_L\).  This is membership through an embedded uniform-half-disc subclass, not a claim of design-wide orthogonality.

The distribution of \((T,Q_\Delta(R))\) is common to the two signs.  In the treated first bin, the difference of their joint success masses is
\[
 \frac1\pi\int_{\Theta_1}\int_0^\Delta
 2c_L\Delta\phi(r/\Delta)r\,dr\,d\vartheta
 =2c_L\Delta^3\int_0^1u\phi(u)\,du=0.
\]
Outside that bin the perturbation vanishes, and all control cells are unchanged.  Because \(Y\) is binary in these two constructed laws, equality of a released cell's total mass and success mass also gives equality of its failure mass.  Therefore the complete one-unit laws of \(Z_\Delta=(Y,T,Q_\Delta(R))\) coincide.  Finally, \(\mu_t(b)=\theta_t\) and \(\phi(0)=1\) give
\[
 \tau_{P^{\mathrm{alias}}_{+,\Delta}}(b)=c_L\Delta,
 \qquad
 \tau_{P^{\mathrm{alias}}_{-,\Delta}}(b)=-c_L\Delta,
\]
whose difference is \(2c_L\Delta\).